% ============================================
% Johns Hopkins Letter Template
% Assistant Professor Cover Letter – University of Northern Iowa
% ============================================

\documentclass[11pt,letterpaper]{letter}

\usepackage{graphicx}
\usepackage{eso-pic}
\usepackage{geometry}
\usepackage{microtype}
\usepackage[hidelinks]{hyperref}

% ----- Page Setup -----
\geometry{left=25mm,right=25mm,top=25mm,bottom=25mm}

% ----- Logo Placement (foreground, first page only) -----
\newcommand{\PutJHULogoFG}[1]{%
  \AddToShipoutPictureFG*{%
    \AtPageUpperLeft{%
      \hspace{10mm}%
      \raisebox{-38mm}{%
        \includegraphics[width=6cm]{#1}%
      }%
    }%
  }%
}

% ----- Sender Information -----
\signature{Decory Edwards}
\address{Department of Economics \\ Johns Hopkins University \\ Baltimore, MD 21218 \\ \texttt{dedwar65@jh.edu}}

\begin{document}
\PutJHULogoFG{Images/jhulogo.png}

\begin{letter}{Search Committee \\
Department of Economics \\
Wilson College of Business \\
University of Northern Iowa \\
Cedar Falls, IA 50614 \\ USA}

\opening{Dear Members of the Search Committee,}

I am writing to express my interest in the tenure-track Assistant Professor of Economics (Macroeconomics) position at the University of Northern Iowa, starting August 2026. I am a Ph.D. candidate in Economics at Johns Hopkins University, advised by Professors Christopher D. Carroll, Nicholas W. Papageorge, and Stelios Fourakis, and I expect to receive my degree in May 2026. My research fields are macroeconomics, wealth and income dynamics, and computational economics.

My research agenda is unified by a focus on understanding the determinants of wealth inequality and the mechanisms through which heterogeneity in returns to wealth shapes aggregate outcomes. In my job-market paper, I estimate the degree of heterogeneity in individual portfolio returns using micro-data from the Survey of Consumer Finances and simulate a structural model in which households face idiosyncratic return risk. I show that allowing for persistent heterogeneity in returns substantially improves the model’s ability to match the empirical distribution of wealth, offering a tractable way to connect micro-level return dispersion to macroeconomic inequality.

In a second project, I use the Health and Retirement Study to study the relationship between interpersonal trust and portfolio performance. Building on the literature that finds a hump-shaped relationship between trust and income, I show that a similar relationship holds for individual returns to wealth measured in the HRS data, highlighting a behavioral channel through which social attitudes shape saving and investment outcomes. This work also provides new estimates of the persistent component of returns to net worth, which motivate the calibration of heterogeneity in my structural model.

The Department of Economics at UNI and the Wilson College of Business emphasize excellence in teaching, the teacher-scholar model, and applied macroeconomics, quantitative methods, and business-economic linkages. The position’s regular teaching of both principles and intermediate macroeconomics, along with opportunities for money \& banking, statistics, or business analytics courses, aligns strongly with my teaching experience and interests. I am particularly excited by the opportunity to bridge theory and practice, bring a global perspective to macroeconomics, and help students understand the role of macroeconomic analysis in business and society.

\textbf{Teaching.} My teaching experience centers on undergraduate macroeconomics and an upper‐level macro policy seminar, where I have led weekly discussion sections and held regular office hours for classes ranging from 16 to 40 students. My approach is student‐centered: I aim to help students “convince themselves” that they have moved from confusion to understanding by holding structured office‐hour coaching that builds trust and confidence. At UNI I am prepared to teach \textit{Intermediate Macroeconomic Theory}, \textit{Principles of Macroeconomics}, \textit{Money \& Banking}, and quantitative methods for economics majors. 

I believe I would be a strong fit because my computational and empirical toolkit allows me to (i) produce robust, data-backed estimates of returns and distributional outcomes, (ii) build and validate models that speak to macroeconomic issues and business-economic questions, and (iii) collaborate with faculty and students to present insights in concise, actionable formats for diverse audiences. I am enthusiastic about joining a collegial, student-focused department committed to teaching excellence and engaged scholarship.

I would be honored to contribute to the University of Northern Iowa’s teaching, scholarship, and service missions and to mentor students pursuing quantitative careers in economics and business. Thank you for your consideration; I welcome the opportunity to discuss my fit with the department at your convenience.

\closing{Sincerely,}

\end{letter}
\end{document}
